\section{Valuable Research Directions}

\subsection{Generalizing Isoperimetric Inequality}

Based on the fact that mean curvature flow constitutes the gradient flow associated with surface area, it is not difficult to deduce that its lower-dimensional counterpart—curve shortening flow—serves as the gradient flow associated with curve length. Drawing an analogy to how mean curvature flow can be utilized to address the minimal surface problem, we are naturally led to wonder: could methods derived from curve shortening flow potentially offer a novel proof of the isoperimetric inequality? The answer is affirmative; in fact, the mathematician Grayson reached a similar conclusion as early as 1983\cite{gage1983isoperimetric}. Our task, therefore, is to clearly distill this proof from his paper and the findings derived from other relevant literature. Specifically, the following references would be helpful\cite{grayson1989shortening,gage1984curve,gage1986heat,grayson1987heat}.

Another thing is that we might be able to apply the method of mean curvature flow to generalize the isoperimetric inequality to higher-dimensional manifolds. For example here is a conjecture
\begin{conjecture}\label{conj:4-1-1}
    Let $M^n$ be a compact, boundaryless hypersurface embedded in $n$-dimensional Euclidean space, enclosing a bounded region $\Omega$. Denoting the $(n+1)$-dimensional volume of this region as $V(\Omega)$ and the $n$-dimensional surface area of $M$ as $A(M)$. Then for some given positive real number $A_0$ there should be a hypersurface that maximizes $V(\Omega)$ under a fixed surface area $A(M)=A_0$.
    Furthermore, for any dimension $n$ there should be a constant $c_n$ that solely depents on $n$ and also makes the following inequality holds.
    \begin{equation}
        \frac{A^{n+1}(M)}{V^n(\Omega)} \ge c_n
    \end{equation}
\end{conjecture}

\subsection{Review the Proof of the Multiplicity-One Conjecture and Try to Prove The Smale Conjecture}